\documentclass[12pt,a4paper]{article}

\usepackage[a4paper,margin=1in]{geometry}
\usepackage{graphicx}
\usepackage{booktabs}
\usepackage{amsmath,amssymb}
\usepackage{float}
\usepackage{hyperref}
\usepackage{xcolor}
\usepackage{listings}
\usepackage{enumitem}
\usepackage{fancyhdr}
\usepackage{longtable}
\usepackage{array}

% Scale a table down to the text width only if it would otherwise overflow.
\newsavebox{\fitbox}
\newenvironment{fitwidth}
  {\begin{lrbox}{\fitbox}}
  {\end{lrbox}\ifdim\wd\fitbox>\linewidth\resizebox{\linewidth}{!}{\usebox{\fitbox}}\else\usebox{\fitbox}\fi}
% Let long code identifiers (e.g. load\_breast\_cancer) wrap at underscores instead of running into the margin.
\DeclareRobustCommand{\_}{\textunderscore\allowbreak}
\setlength{\emergencystretch}{3em}

\hypersetup{colorlinks=true,linkcolor=blue,urlcolor=blue}

\pagestyle{fancy}
\fancyhf{}
\lhead{ICS1512}
\rhead{Machine Learning Algorithms Laboratory}
\cfoot{\thepage}

\lstset{
language=Python,
basicstyle=\ttfamily\small,
numbers=left,
frame=single,
breaklines=true,
showstringspaces=false,
keywordstyle=\color{blue},
commentstyle=\color{gray},
stringstyle=\color{purple}
}

\begin{document}

\noindent\begin{tabular}{|p{4cm}|p{\dimexpr\linewidth-4cm-4\tabcolsep-3\arrayrulewidth\relax}|}
\hline
Experiment No. & 3 \\ \hline
Experiment Title & Regression Analysis using Linear and Regularized Models
\footnote{\textbf{GitHub Repository: }\url{https://github.com/Nightingales21/ICS1512-Machine-Learning-Labratory}}\\ \hline
Student Name & Nitin \\ \hline
Register Number & 312224700142 \\ \hline
Faculty & Dr. Poreddy Ajay Kumar Reddy \\ \hline
Submission Date & \today \\ \hline
\end{tabular}

\section{Aim and Objective}
To implement linear and regularized regression models (Linear, Ridge, Lasso,
and Elastic Net) for predicting a continuous target variable -- the
sanctioned loan amount -- evaluate their performance using multiple
regression metrics, visualize model behavior, and analyze overfitting,
underfitting, and bias--variance characteristics.

\section{Dataset Description}
\begin{longtable}{|p{6cm}|p{8cm}|}
\hline
Dataset Name & Loan Amount Prediction (Analytics Vidhya / Kaggle ``Predict
Loan Amount Data'') \\ \hline
Dataset Source & Kaggle, originally hosted via the Analytics Vidhya Loan
Prediction Practice Problem \\ \hline
Number of Samples (raw) & 614 loan applications \\ \hline
Number of Samples (used) & 592 (after dropping 22 rows with missing target) \\ \hline
Number of Features & 11 predictor features (4 numeric, 7 categorical) after
dropping the identifier column \\ \hline
Target Variable & \texttt{LoanAmount} (continuous, in thousands) \\ \hline
Missing Values & Present in \texttt{Gender} (13), \texttt{Married} (3),
\texttt{Dependents} (15), \texttt{Self\_Employed} (32), \texttt{LoanAmount}
(22), \texttt{Loan\_Amount\_Term} (14), \texttt{Credit\_History} (50) \\ \hline
Train--Test Split & 80\% / 20\%, \texttt{random\_state=42} \\ \hline
\end{longtable}

Feature list: \texttt{Gender}, \texttt{Married}, \texttt{Dependents},
\texttt{Education}, \texttt{Self\_Employed}, \texttt{ApplicantIncome},
\texttt{CoapplicantIncome}, \texttt{Loan\_Amount\_Term},
\texttt{Credit\_History}, \texttt{Property\_Area}, \texttt{Loan\_Status}.
The \texttt{Loan\_ID} identifier column was dropped as it carries no
predictive information.

\section{Preprocessing Steps}
\begin{itemize}
\item \textbf{Target rows with missing target:} The 22 rows with a missing
\texttt{LoanAmount} were dropped, since the target cannot be imputed without
introducing circular bias into a regression task.
\item \textbf{Missing value handling (features):} Categorical features
(\texttt{Gender}, \texttt{Married}, \texttt{Dependents}, \texttt{Self\_Employed})
were imputed with their mode; numeric features (\texttt{Loan\_Amount\_Term},
\texttt{Credit\_History}) were imputed with their median.
\item \textbf{Encoding:} \texttt{Dependents} (ordinal, with ``3+'' mapped to
3) was kept numeric; the remaining categorical features (\texttt{Gender},
\texttt{Married}, \texttt{Education}, \texttt{Self\_Employed},
\texttt{Property\_Area}, \texttt{Loan\_Status}) were one-hot encoded with
\texttt{drop\_first=True} to avoid the dummy-variable trap.
\item \textbf{Feature scaling:} \texttt{StandardScaler} was fit on the
training split only and applied to both train and test features before
every model (Linear, Ridge, Lasso, Elastic Net), since regularization
penalties are scale-dependent.
\item \textbf{Train--test split:} \texttt{train\_test\_split} with
\texttt{test\_size=0.2}, \texttt{random\_state=42}.
\end{itemize}

\section{Exploratory Data Analysis}

The reusable function \texttt{generate\_eda\_summary()} from
\texttt{ml\_lab\_utils.py} (developed in Experiment 1, unmodified) was used
to produce the consolidated 12-panel EDA figure below.

\begin{figure}[H]
\centering
\includegraphics[width=\linewidth]{figures/eda_loan.png}
\caption{Consolidated 12-panel EDA summary for the Loan Amount Prediction dataset.}
\label{fig:eda}
\end{figure}

\textbf{Inference:} After imputation there are no remaining missing values.
\texttt{Applicant\allowbreak Income} is the highest-variance numeric feature by a wide
margin and is strongly right-skewed (Panels 6, 9, 10, 12), with a small
number of very high-income applicants visible as outliers in the box plot
(Panel 7) and as the far-right point in the Q--Q plot's upper tail (Panel
10). \texttt{LoanAmount} itself (Panel 4) is also right-skewed, which is
consistent with the residual patterns observed later in Section~6.
\texttt{ApplicantIncome} and \texttt{CoapplicantIncome} show a weak negative
correlation (Panel 5, $-0.11$), suggesting the two income sources are
largely independent (i.e., loan applications are not being systematically
structured to trade off one income source against the other).

\subsection*{Additional Required Visualizations}

\begin{figure}[H]
\centering
\includegraphics[width=0.7\linewidth]{figures/target_distribution.png}
\caption{Target variable (\texttt{LoanAmount}) distribution.}
\label{fig:target}
\end{figure}
\textbf{Inference:} The target is unimodal and right-skewed, with most
sanctioned loan amounts falling between 100 and 170 (thousand), and a long
tail of high-value loans extending past 400 -- these tail cases are the
hardest for a linear model to predict accurately, as confirmed in the
predicted-vs-actual plots (Figure~\ref{fig:pred_actual}).

\begin{figure}[H]
\centering
\includegraphics[width=\linewidth]{figures/feature_vs_target_scatter.png}
\caption{Applicant income and coapplicant income vs. loan amount.}
\label{fig:feat_target}
\end{figure}
\textbf{Inference:} Both income features show a weak positive, high-variance
relationship with loan amount -- higher income allows for (but does not
guarantee) a larger loan, consistent with the low $R^2$ values obtained
later in Section~6. Applicant income shows a slightly stronger and more
visible upward trend than coapplicant income.

\section{Implementation Details}

\textbf{Linear Regression} models the relationship between input features
and the continuous target as $\hat{y} = \boldsymbol\beta^\top\mathbf{x} +
\beta_0$, fit by minimizing squared error. It is simple, interpretable, and
serves as this experiment's baseline.

\textbf{Ridge Regression} adds an $L_2$ penalty
$\alpha\sum_j\beta_j^2$ to the ordinary least-squares objective, shrinking
coefficient magnitudes toward zero (but not exactly to zero) and reducing
variance at the cost of a small increase in bias.

\textbf{Lasso Regression} adds an $L_1$ penalty $\alpha\sum_j|\beta_j|$,
which can shrink coefficients exactly to zero, performing automatic feature
selection.

\textbf{Elastic Net} combines both penalties,
$\alpha\left(\text{l1\_ratio}\sum_j|\beta_j| + \tfrac{1-\text{l1\_ratio}}{2}\sum_j\beta_j^2\right)$,
balancing Ridge's stability with Lasso's sparsity.

\begin{lstlisting}[language=Python, caption={Reusable regression training call, from \texttt{ml\_lab\_utils.py} (Experiment 1) applied to the baseline model.}]
baseline_models = {"Linear Regression": LinearRegression()}
baseline_results_df, baseline_fitted = train_evaluate_regression(
    baseline_models, X_train_scaled, X_test_scaled, y_train, y_test,
    scale=False)
\end{lstlisting}

\begin{lstlisting}[language=Python, caption={Hyperparameter search space and GridSearchCV/RandomizedSearchCV tuning, from \texttt{experiment3\_loan\_regression.py}.}]
search_spaces = {
    "Ridge Regression": (Ridge(random_state=42),
                          {"alpha": [0.01, 0.1, 1, 10, 100]}),
    "Lasso Regression": (Lasso(random_state=42, max_iter=10000),
                          {"alpha": [0.001, 0.01, 0.1, 1, 10]}),
    "Elastic Net Regression": (ElasticNet(random_state=42, max_iter=10000),
                                {"alpha": [0.01, 0.1, 1, 10],
                                 "l1_ratio": [0.2, 0.5, 0.8]}),
}
gs = GridSearchCV(estimator, grid, cv=KFold(5, shuffle=True, random_state=42),
                   scoring="r2", n_jobs=-1)
gs.fit(X_train_scaled, y_train)
\end{lstlisting}

\section{Experimental Setup}
\begin{center}
\begin{minipage}{\linewidth}\centering
\begin{fitwidth}
\begin{tabular}{ll}
\toprule
Python Version & 3.12 \\
Scikit-Learn Version & 1.8.0 \\
IDE & Jupyter Notebook / VS Code \\
Operating System & Ubuntu 24.04 LTS \\
Hardware & Standard lab workstation (CPU only) \\
\bottomrule
\end{tabular}
\end{fitwidth}
\end{minipage}
\end{center}

\vspace{0.3cm}
\textbf{Environment activation command (as mandated by the lab manual):}
\begin{lstlisting}[language=bash]
source /opt/anaconda3/bin/activate anaconda-navigation
\end{lstlisting}

\section{Hyperparameter Tuning Results}

\begin{center}
\begin{minipage}{\linewidth}\centering
\textbf{Table 1: Hyperparameter Tuning Summary}\par\smallskip
\begin{fitwidth}
\begin{tabular}{lllc}
\toprule
Model & Search Method & Best Parameters & Best CV $R^2$ \\
\midrule
Ridge Regression       & Grid / Random & $\alpha=100$                        & 0.3714 \\
Lasso Regression       & Grid / Random & $\alpha=10$                         & 0.3311 \\
Elastic Net Regression & Grid / Random & $\alpha=1$, l1\_ratio$=0.8$         & 0.3707 \\
\bottomrule
\end{tabular}
\end{fitwidth}
\end{minipage}
\end{center}

\vspace{0.2cm}
\noindent GridSearchCV and RandomizedSearchCV selected \emph{identical}
best hyperparameters and CV $R^2$ for all three regularized models on this
small search space, with RandomizedSearchCV requiring modestly fewer fit
calls (Ridge: 5 vs.\ 5, Lasso: 5 vs.\ 5, Elastic Net: 10 vs.\ 12), since the
grids here are small enough that random sampling covers most or all of the
space.

\section{Performance Tables}

\begin{center}
\begin{minipage}{\linewidth}\centering
\textbf{Table 2: Cross-Validation Performance (K = 5, on training data)}\par\smallskip
\begin{fitwidth}
\begin{tabular}{lcccc}
\toprule
Model & MAE & MSE & RMSE & $R^2$ \\
\midrule
Linear Regression        & 43.11 & 5508.55 & 74.22 & 0.3109 \\
Ridge Regression         & 44.08 & 4935.28 & 70.25 & \textbf{0.3714} \\
Lasso Regression         & 46.23 & 5281.51 & 72.67 & 0.3311 \\
Elastic Net Regression   & 43.88 & 4944.29 & 70.32 & 0.3707 \\
\bottomrule
\end{tabular}
\end{fitwidth}
\end{minipage}
\end{center}

\vspace{0.3cm}
\begin{center}
\begin{minipage}{\linewidth}\centering
\textbf{Table 3: Test Set Performance}\par\smallskip
\begin{fitwidth}
\begin{tabular}{lcccc}
\toprule
Model & MAE & MSE & RMSE & $R^2$ \\
\midrule
Linear Regression       & 36.88 & 3335.43 & 57.75 & 0.0952 \\
Ridge Regression        & 37.55 & 3159.55 & 56.21 & 0.1429 \\
Lasso Regression        & 37.80 & 3176.80 & 56.36 & 0.1382 \\
Elastic Net Regression  & 37.43 & 3151.21 & 56.14 & \textbf{0.1452} \\
\bottomrule
\end{tabular}
\end{fitwidth}
\end{minipage}
\end{center}

\vspace{0.2cm}
\noindent All four models train in well under a millisecond on this
dataset size (592 samples, 12 features); training time is therefore not a
practically distinguishing factor here, and all further discussion focuses
on predictive performance.

\subsection*{Test Set Performance Comparison}
\begin{figure}[H]
\centering
\includegraphics[width=\linewidth]{figures/predicted_vs_actual.png}
\caption{Predicted vs. actual loan amount for all four models on the held-out test set.}
\label{fig:pred_actual}
\end{figure}
\textbf{Inference:} All four models show the same qualitative failure mode:
they compress predictions toward the mean and systematically under-predict
the small number of very large loan amounts (the points far to the right,
well below the dashed $y=x$ line). Ridge and Elastic Net show slightly
tighter clustering around the diagonal for mid-range loan amounts than plain
Linear Regression, consistent with their higher test $R^2$.

\begin{figure}[H]
\centering
\includegraphics[width=\linewidth]{figures/residual_plots.png}
\caption{Residual plots (residual = actual $-$ predicted) for all four models.}
\label{fig:residuals}
\end{figure}
\textbf{Inference:} Residuals are not randomly scattered around zero -- there
is a clear funnel/fan pattern (heteroscedasticity), with residual variance
increasing for larger predicted loan amounts. This, combined with the
positive residuals concentrated at higher predicted values, confirms that a
plain linear model under-fits the (right-skewed, likely non-linear)
relationship between income and loan amount at the high end of the
distribution -- a log-transform of the target would be a natural next step
to address this in future work.

\begin{figure}[H]
\centering
\includegraphics[width=0.75\linewidth]{figures/coefficient_comparison.png}
\caption{Coefficient comparison across Linear, Ridge, Lasso, and Elastic Net (top 8 features by summed magnitude).}
\label{fig:coef}
\end{figure}
\begin{center}
\begin{minipage}{\linewidth}\centering
\textbf{Table 4: Coefficient Comparison (top 3 features)}\par\smallskip
\begin{fitwidth}
\begin{tabular}{lcccc}
\toprule
Feature & Linear & Ridge & Lasso & Elastic Net \\
\midrule
ApplicantIncome     & 53.95 & 43.83 & 44.98 & 43.72 \\
CoapplicantIncome   & 21.97 & 17.00 & 11.79 & 16.63 \\
Married\_Yes        & 9.72  & 8.08  & 1.74  & 7.64  \\
\bottomrule
\end{tabular}
\end{fitwidth}
\end{minipage}
\end{center}
\textbf{Inference:} \texttt{ApplicantIncome} is the dominant predictor across
every model. Ridge shrinks all coefficients smoothly toward zero without
eliminating any of them. Lasso is far more aggressive: it drives
\texttt{Dependents}, \texttt{Loan\_Amount\_Term}, \texttt{Credit\_History},
\texttt{Gender\_Male}, \texttt{Self\_Employed\_Yes}, and three
\texttt{Property\_Area}/\texttt{Loan\_Status} dummies to \emph{exactly}
zero, keeping only \texttt{ApplicantIncome}, \texttt{CoapplicantIncome}, and
\texttt{Married\_Yes} with non-trivial weight. Elastic Net sits between the
two, shrinking coefficients but retaining more non-zero features than Lasso.

\section{Overfitting and Underfitting Analysis}

\begin{figure}[H]
\centering
\includegraphics[width=0.8\linewidth]{figures/train_vs_validation_error.png}
\caption{Training RMSE vs. 5-fold validation RMSE for all four models.}
\label{fig:train_val}
\end{figure}

\textbf{Training vs. validation error:} Plain Linear Regression has the
lowest training RMSE (66.09) but the \emph{largest} training--validation gap
(66.09 $\to$ 73.52, a gap of 7.43), the classic signature of a model fitting
noise in the training data (mild overfitting relative to the regularized
alternatives). Ridge (67.60 $\to$ 69.94, gap 2.35) and Elastic Net (67.34
$\to$ 70.02, gap 2.69) both substantially narrow this generalization gap at
the cost of a small increase in training error -- exactly the bias-variance
trade-off regularization is designed to make. Lasso narrows the gap less
well here (69.99 $\to$ 72.30, gap 2.31 in RMSE terms but from a higher
starting point), consistent with it having discarded several features
(\texttt{Credit\_History} in particular) that do carry real, if modest,
predictive signal.

\textbf{Effect of regularization strength:} Figure~\ref{fig:ridge_path} shows
Ridge's regularization path directly: as $\alpha$ increases from 0.01 to
100, training RMSE rises only slightly (66.09 $\to$ 67.60) while validation
RMSE falls substantially (73.52 $\to$ 69.94). This is the textbook
regularization effect -- controlled bias increase buys a larger variance
reduction, improving generalization, up to the point where too much
regularization would eventually push training error up faster than
validation error falls (not yet reached within the tested $\alpha$ range
here, since $\alpha=100$ was selected as the CV-optimal value by
GridSearchCV).

\begin{figure}[H]
\centering
\includegraphics[width=0.7\linewidth]{figures/ridge_alpha_path.png}
\caption{Ridge: training vs. validation RMSE as a function of $\alpha$.}
\label{fig:ridge_path}
\end{figure}

\textbf{Improvement in generalization after tuning:} Cross-validated $R^2$
improved from 0.3109 (untuned Linear Regression) to 0.3714 (tuned Ridge) --
a meaningful relative improvement of roughly 19\%, achieved purely by
adding an appropriately-tuned $L_2$ penalty with no change to the feature
set.

\section{Bias--Variance Analysis}

\textbf{Bias behavior of Linear Regression:} Ordinary Linear Regression has
the lowest bias among the four models (it best fits the training data, as
shown by its lowest training RMSE), but this low bias comes at the cost of
higher variance -- it is the most sensitive to the particular training fold,
producing the largest train--validation gap observed in this experiment.

\textbf{Variance reduction using Ridge and Elastic Net:} Both Ridge and
Elastic Net shrink coefficient magnitudes (Figure~\ref{fig:coef}), which
directly reduces the model's sensitivity to small changes in the training
data (variance). This is empirically confirmed by their smaller
train--validation RMSE gaps (2.35 and 2.69 respectively) compared to plain
Linear Regression's gap of 7.43, and by their higher cross-validated $R^2$
(0.3714 and 0.3707 vs. 0.3109).

\textbf{Feature sparsity effect in Lasso:} Figure~\ref{fig:lasso_sparsity}
shows the number of non-zero coefficients falling from 12 (at
$\alpha=0.001$) to just 3 (at $\alpha=10$, the CV-selected optimum) as
regularization strength increases. This sparsity is a form of automatic
feature selection: Lasso concludes that only \texttt{ApplicantIncome},
\texttt{CoapplicantIncome}, and \texttt{Married\_Yes} carry enough signal to
justify a non-zero coefficient at its chosen $\alpha$, trading a small amount
of additional bias (its test $R^2$ of 0.1382 is the lowest of the three
regularized models) for a much simpler, more interpretable model.

\begin{figure}[H]
\centering
\includegraphics[width=0.7\linewidth]{figures/lasso_sparsity.png}
\caption{Lasso: number of non-zero coefficients vs. regularization strength $\alpha$.}
\label{fig:lasso_sparsity}
\end{figure}

\section{Observations and Conclusion}
All four models achieve broadly similar test-set performance (test $R^2$
between 0.095 and 0.145), reflecting a genuinely difficult regression
target: loan amount sanctioned depends on underwriting factors (e.g.,
detailed credit assessment, property valuation) not captured in this
feature set, and the target itself is right-skewed with influential
high-value outliers. Among the four models, \textbf{Elastic Net Regression}
achieved the best test-set $R^2$ (0.1452) and joint-best RMSE (56.14),
narrowly ahead of Ridge (0.1429, 56.21); both clearly outperform plain
Linear Regression (0.0952, 57.75) and Lasso (0.1382, 56.36). The optimal
hyperparameters -- Ridge $\alpha=100$, Lasso $\alpha=10$, Elastic Net
$\alpha=1$ with l1\_ratio$=0.8$ -- were each confirmed identically by both
GridSearchCV and RandomizedSearchCV, giving confidence that these are true
optima within the specified search space rather than artifacts of a
particular search strategy. In terms of the accuracy-versus-complexity
trade-off: Elastic Net offers the best accuracy but retains most features
with shrunk (not zeroed) weights; Lasso sacrifices a small amount of
accuracy for a much simpler 3-feature model, which may be preferable in a
setting where model interpretability or deployment simplicity outweighs the
last few percentage points of $R^2$. For this dataset, given the modest
improvement in test $R^2$ (0.143 vs 0.145) between Ridge and Elastic Net,
either would be a reasonable choice; the sparser Lasso model is the more
attractive option if explainability to loan officers is a priority, despite
its slightly lower accuracy.

\section{References}
\begin{enumerate}
\item Scikit-learn Developers, ``Linear Models,'' scikit-learn
documentation. [Online]. Available:
\url{https://scikit-learn.org/stable/modules/linear_model.html}
\item Scikit-learn Developers, ``Tuning the hyper-parameters of an
estimator,'' scikit-learn documentation. [Online]. Available:
\url{https://scikit-learn.org/stable/modules/grid_search.html}
\item Analytics Vidhya / Kaggle, ``Loan Prediction Problem Dataset.''
[Online]. Available:
\url{https://www.kaggle.com/datasets/altruistdelhite04/loan-prediction-problem-dataset}
\item F. Pedregosa et al., ``Scikit-learn: Machine Learning in Python,''
\emph{Journal of Machine Learning Research}, vol. 12, pp. 2825--2830, 2011.
\end{enumerate}

\end{document}
